\section{Commuting orbit rows and joint moments}
\label{sec:rows}

Throughout this section, \(A\), \(G\), and \(\alpha\) satisfy the
hypotheses of Theorem~\ref{thm:main}. Only countability of \(G\) is
needed for the constructions in this section.

\begin{proposition}[Commuting orbit rows]
\label{prop:rows}
There are unital embeddings
\(\iota_i\colon\Zalg\to A^\omega\), \(i\geq 1\), for which the unital
orbit algebras
\[
 C_i=C^*\bigl(\bar\alpha_g(\iota_i(\Zalg)):g\in G\bigr)
\]
commute with \(A\) and commute pairwise. In particular,
\[
 C=C^*\Bigl(\bigcup_{i\geq 1}C_i\Bigr)
 \subseteq A^\omega\cap A'
\]
is unital, separable, and \(\bar\alpha\)-invariant, and
\(\T(C)\neq\varnothing\).

Moreover, there are self-adjoint elements \(x_{i,g}\in C_i\),
\(i\geq 1\), \(g\in G\), and a constant \(M<\infty\), such that
\begin{equation}
\label{eq:row-normalization}
 \|x_{i,g}\|\leq M,\qquad
 \varphi(x_{i,g})=0,\qquad
 \varphi(x_{i,g}^2)=1
 \quad\bigl(\varphi\in\T(C)\bigr),
\end{equation}
and
\begin{equation}
\label{eq:row-equivariance}
 \bar\alpha_h(x_{i,g})=x_{i,hg}
 \qquad(i\geq1,\ g,h\in G).
\end{equation}
\end{proposition}

\begin{proof}
The Jiang--Su algebra is simple, unital, uniquely tracial, and
isomorphic to its infinite tensor power; see
\cite[Example~1.14(iv) and Corollary~1.11]{TomsWinter2007}.
Write \(\rho\) for its unique tracial state. Since \(A\) is
\(\Zalg\)-stable, fix an identification
\[
 A\cong A\otimes\bigotimes_{l=1}^{\infty}\Zalg.
\]
We use the increasing finite tensor stages in this identification.

We first show that every norm-separable \(C^*\)-subalgebra
\(D\subseteq A^\omega\) admits a unital embedding
\begin{equation}
\label{eq:relative-tail-embedding}
 \Zalg\longrightarrow A^\omega\cap D'.
\end{equation}
Choose a norm-dense sequence \((d_s)_{s\geq1}\) in \(D\), and choose
bounded representative sequences
\(d_s=[(d_{s,j})_{j\geq1}]\).
For each \(j\), choose a finite tensor stage containing elements
within \(1/j\) in norm of
\(d_{1,j},\ldots,d_{j,j}\). Let
\(\theta_j\colon\Zalg\to A\) be the embedding into a tensor factor
strictly after this stage. For \(j\geq s\) and \(z\in\Zalg\),
\[
 \|[\theta_j(z),d_{s,j}]\|\leq \frac{2\|z\|}{j}.
\]
The coordinatewise homomorphisms therefore induce a unital
\(*\)-homomorphism \(z\mapsto[(\theta_j(z))_j]\) into
\(A^\omega\cap D'\). It is injective because \(\Zalg\) is simple
and the unit of \(A^\omega\) is nonzero. This proves
Equation~\eqref{eq:relative-tail-embedding}.

Set \(D_0=A\). Inductively, suppose that \(D_{i-1}\) is a separable,
unital, \(\bar\alpha\)-invariant subalgebra of \(A^\omega\) containing
\(A\) and the preceding orbit algebras. Apply
Equation~\eqref{eq:relative-tail-embedding} to obtain
\(\iota_i(\Zalg)\subseteq D_{i-1}'\), define \(C_i\) as in the
statement, and put
\[
 D_i=C^*(D_{i-1},C_i).
\]
For \(b\in\iota_i(\Zalg)\), \(d\in D_{i-1}\), and \(g\in G\),
\[
 [\bar\alpha_g(b),d]
 =\bar\alpha_g\bigl([b,\bar\alpha_{g^{-1}}(d)]\bigr)=0,
\]
since \(D_{i-1}\) is invariant. Thus \(C_i\) commutes with
\(D_{i-1}\). Countability of \(G\) preserves separability at each
stage. This gives the required algebras and their commutation
properties. A limit trace on \(A^\omega\) restricts to a tracial
state on \(C\), so \(\T(C)\) is nonempty.

The trace \(\rho\) is faithful by simplicity of \(\Zalg\).
Choose a non-scalar self-adjoint element \(c\in\Zalg\). Then
\(\rho((c-\rho(c)1)^2)>0\), so rescaling \(c-\rho(c)1\) gives
a self-adjoint \(b\in\Zalg\) with
\(\rho(b)=0\) and \(\rho(b^2)=1\). Set
\[
 x_{i,g}=\bar\alpha_g(\iota_i(b)),\qquad M=\|b\|.
\]
Every \(\varphi\in\T(C)\) restricts along each translated unital
embedding of \(\Zalg\) to \(\rho\). This proves
Equation~\eqref{eq:row-normalization}, and
Equation~\eqref{eq:row-equivariance} follows from the action law.
\end{proof}

The following independence statement is intrinsic to the algebra
generated by the rows. In particular, it applies to traces of \(C\)
that do not extend to the ambient ultrapower.

\begin{lemma}[Independence at extreme traces]
\label{lem:factor-independence}
Let \(B\) be a unital \(C^*\)-algebra generated by pairwise commuting
unital \(C^*\)-subalgebras \(B_i\), \(i\geq1\). If
\(\varphi\in\partial_{\mathrm e}\T(B)\), then
\[
 \varphi(b_1\cdots b_s)=\prod_{j=1}^s\varphi(b_j)
\]
whenever \(b_j\in B_{i_j}\) and the indices \(i_1,\ldots,i_s\)
are distinct.
\end{lemma}

\begin{proof}
Let \((H_\varphi,\pi,\Omega)\) be the tracial GNS representation
of \(\varphi\), and put
\[
 \mathcal M=\pi(B)'',\qquad
 \mathcal M_i=\pi(B_i)''.
\]
We first record the relevant properties of this representation,
allowing \(\varphi\) to have a nonzero kernel on \(B\).
On the dense subspace of vectors \([a]\), \(a\in B\), right
multiplication \(R_b[a]=[ab]\) is bounded with norm at most
\(\|b\|\), by traciality. It commutes with \(\pi(B)\), while
the vectors \(R_b\Omega=[b]\) are dense. If
\(T\in\mathcal M\) satisfies \(T\Omega=0\), then
\[
 TR_b\Omega=R_bT\Omega=0\qquad(b\in B),
\]
so \(T=0\). Thus \(\Omega\) is separating for \(\mathcal M\).
The normal vector state
\(\tau(T)=\langle T\Omega,\Omega\rangle\) is consequently
faithful: for \(T\geq0\), \(\tau(T)=0\) implies
\(T^{1/2}\Omega=0\). Its traciality extends from \(\pi(B)\)
to \(\mathcal M\) by separate ultraweak continuity. Hence
\(\mathcal M\) is a finite von Neumann algebra with faithful
normal tracial state \(\tau\).

Extremality of \(\varphi\) implies that \(\mathcal M\) is a factor.
Indeed, a nontrivial central projection \(p\in\mathcal M\) would
give the convex decomposition
\[
 \varphi(a)
 =\tau(p)\frac{\tau(p\pi(a))}{\tau(p)}
  +\tau(1-p)\frac{\tau((1-p)\pi(a))}{\tau(1-p)}.
\]
The two normalized functionals are tracial states on \(B\). They are distinct:
if, for instance, the first equalled \(\varphi\), equality of
the corresponding normal functionals would extend from
\(\pi(B)\) to \(\mathcal M\), and evaluation at \(p\) would
give \(1=\tau(p)\). This contradicts
\(0<\tau(p)<1\).

The algebras \(\mathcal M_i\) commute and generate \(\mathcal M\).
Every element of the center of \(\mathcal M_i\) therefore
commutes with all of \(\mathcal M\), so
\[
 Z(\mathcal M_i)\subseteq Z(\mathcal M)=\C1.
\]
The restriction of \(\tau\) makes each \(\mathcal M_i\) a finite
factor with a faithful normal tracial state.

Fix \(i\), and let \(y\geq0\) belong to the von Neumann algebra
generated by the other \(\mathcal M_j\)'s. It commutes with
\(\mathcal M_i\). The functional
\[
 \psi_y(x)=\tau(y^{1/2}xy^{1/2})=\tau(xy),
 \qquad x\in\mathcal M_i,
\]
is normal and positive: sandwiching by \(y^{1/2}\) preserves
suprema of bounded increasing nets of positive operators.
Commutation of \(y\) with \(\mathcal M_i\) makes it tracial,
and \(0\leq\psi_y\leq\|y\|\tau\).
The Radon--Nikodym Theorem for normal functionals on a finite
von Neumann algebra therefore gives a density
\(h\in\mathcal M_i\), \(0\leq h\leq\|y\|1\), with
\(\psi_y(x)=\tau(hx)\). Invariance of \(\psi_y\) under unitary
conjugation forces \(h\) to commute with every unitary in
\(\mathcal M_i\). Hence \(h\) is central and thus scalar.
Evaluating at the unit identifies that scalar as \(\tau(y)\).
It follows that
\[
 \tau(xy)=\tau(x)\tau(y).
\]
Linearity extends this identity to arbitrary \(y\) in the
algebra generated by the other rows. Applying it successively
proves the required product formula.
\end{proof}

We now fix the algebras and elements from
Proposition~\ref{prop:rows}. For \(N\geq1\) and \(g,h\in G\), set
\begin{equation}
\label{eq:empirical-variables}
 z_{N,g}=\frac1{\sqrt N}\sum_{i=1}^N x_{i,g},
 \qquad
 k_{N;g,h}=\frac1N\sum_{i=1}^N
 \frac{x_{i,g}x_{i,h}+x_{i,h}x_{i,g}}2.
\end{equation}
The last two subscripts of \(k_{N;g,h}\) are covariance indices.
These elements are self-adjoint, \(k_{N;g,h}=k_{N;h,g}\), and
\begin{equation}
\label{eq:empirical-equivariance}
 \bar\alpha_l(z_{N,g})=z_{N,lg},
 \qquad
 \bar\alpha_l(k_{N;g,h})=k_{N;lg,lh}
 \quad(l,g,h\in G).
\end{equation}

\begin{proposition}[Joint ordered moment limit]
\label{prop:joint-clt}
Let \(J\subseteq G\) be nonempty and finite. Let \(N_j\to\infty\)
be a sequence of positive integers, and let
\(\varphi_j\in\partial_{\mathrm e}\T(C)\). Suppose that
\begin{equation}
\label{eq:covariance-convergence}
 \varphi_j(k_{N_j;g,h})\longrightarrow R_{g,h}
 \qquad(g,h\in J).
\end{equation}
Then \(R=(R_{g,h})_{g,h\in J}\) is a real positive semidefinite
matrix with diagonal entries one.
Choose unit vectors \(v_g\), \(g\in J\), spanning a
finite-dimensional real Euclidean space \(L\), with
\(\langle v_g,v_h\rangle=R_{g,h}\). If \(W\) is a standard
Gaussian vector in \(L\) and
\(\xi_g=\langle v_g,W\rangle\), then
\[
 \bigl((z_{N_j,g})_{g\in J},(k_{N_j;g,h})_{g,h\in J}\bigr)
 \longrightarrow
 \bigl((\xi_g)_{g\in J},(R_{g,h})_{g,h\in J}\bigr)
\]
in all ordered joint moments under \(\varphi_j\).
Equivalently, for every noncommutative \(*\)-polynomial \(P\)
in these self-adjoint variables,
\begin{equation}
\label{eq:joint-polynomial-limit}
 \varphi_j\!\left(
 P\bigl((z_{N_j,g})_g,(k_{N_j;g,h})_{g,h}\bigr)\right)
 \longrightarrow
 \mathbb E\,P\bigl((\xi_g)_g,(R_{g,h})_{g,h}\bigr).
\end{equation}
For every sequence \((N_j,\varphi_j)\) as above, a subsequence
satisfying Equation~\eqref{eq:covariance-convergence} exists.
\end{proposition}

\begin{proof}
For a tracial state \(\varphi\) and self-adjoint \(x,y\),
\[
 \varphi(xy)=\varphi(yx)
 =\overline{\varphi(xy)}.
\]
Thus all the second moments under consideration are real
and symmetric. For real scalars \(a_g\),
\[
 \sum_{g,h\in J}a_ga_h\varphi_j(k_{N_j;g,h})
 =
 \frac1{N_j}\sum_{i=1}^{N_j}
 \varphi_j\!\left(\left(\sum_{g\in J}a_gx_{i,g}\right)^2\right)
 \geq0.
\]
The diagonal entries are one by
Equation~\eqref{eq:row-normalization}. Cauchy--Schwarz bounds every
entry in absolute value by one. Consequently these covariance
matrices lie in a compact set of real symmetric matrices;
their limits have the asserted properties, and the final
subsequence assertion follows. A Gram representation of
the positive semidefinite matrix \(R\) supplies \(L\) and
the vectors \(v_g\).

The empty word is immediate. We first consider a word of
length \(d\geq1\) in the fluctuation
variables. At one term of the given sequence write
\(N=N_j\) and \(\varphi=\varphi_j\). Expanding the sums gives
\begin{equation}
\label{eq:fluctuation-expansion}
 \varphi(z_{N,g_1}\cdots z_{N,g_d})
 =N^{-d/2}
 \sum_{i_1,\ldots,i_d=1}^N
 \varphi(x_{i_1,g_1}\cdots x_{i_d,g_d}).
\end{equation}
Partition the positions \(1,\ldots,d\) according to equality
of their row labels \(i_a\). Since distinct rows commute,
we may group the factors belonging to each row, preserving
their order within that row. Lemma~\ref{lem:factor-independence}
then factors the expectation over the distinct rows.
A singleton block gives zero by
Equation~\eqref{eq:row-normalization}. Each remaining summand has
absolute value at most \(M^d\).

A partition with \(b\) blocks has at most \(N^b\) choices of
distinct row labels. Its contribution to
Equation~\eqref{eq:fluctuation-expansion} is therefore bounded by
\(M^dN^{b-d/2}\). Among partitions without singleton blocks,
\(b\leq\lfloor d/2\rfloor\). It follows that odd moments
tend to zero, and, when \(d=2q\), only partitions into \(q\)
pairs can have a nonzero limit.

Fix such a pairing, writing its pairs as
\((a_l,b_l)\), \(1\leq l\leq q\), with \(a_l<b_l\).
Its contribution is
\[
 N^{-q}
 \sum_{\substack{1\leq i_1,\ldots,i_q\leq N\\
                  i_1,\ldots,i_q\ \mathrm{distinct}}}
 \prod_{l=1}^q
 \varphi(x_{i_l,g_{a_l}}x_{i_l,g_{b_l}}).
\]
Dropping the distinctness condition changes this by
\(O_{d,M}(N^{-1})\): there are at most
\(\binom q2 N^{q-1}\) tuples with a collision, and each
product is bounded by \(M^{2q}\).
The unrestricted sum factors, and traciality identifies
its factors with the empirical covariances:
\[
 \prod_{l=1}^q
 \left(\frac1N\sum_{i=1}^N
 \varphi(x_{i,g_{a_l}}x_{i,g_{b_l}})\right)
 =
 \prod_{l=1}^q\varphi(k_{N;g_{a_l},g_{b_l}}).
\]
All error estimates depend only on \(d\) and \(M\), so they
hold uniformly for the moving traces \(\varphi_j\).
By Equation~\eqref{eq:covariance-convergence}, the limiting even
moment is the sum over pairings of the corresponding
products of entries of \(R\). This is the Gaussian moment:
the moment-generating function of \((\xi_g)_{g\in J}\) is
\[
 \mathbb E\exp\!\left(\sum_{g\in J}t_g\xi_g\right)
 =
 \exp\!\left(\frac12\sum_{g,h\in J}t_gt_hR_{g,h}\right),
\]
whose Taylor expansion gives precisely this pairing formula.

It remains to include the covariance variables in arbitrary
positions of a word. For \(g,h\in J\), define
\[
 s_{i;g,h}=\frac{x_{i,g}x_{i,h}+x_{i,h}x_{i,g}}2,
 \qquad
 y_{j,i;g,h}=s_{i;g,h}-\varphi_j(s_{i;g,h})1.
\]
This is centering separately in each row, at the current
trace. In particular,
\[
 \varphi_j(y_{j,i;g,h})=0,\qquad
 \|y_{j,i;g,h}\|\leq2M^2,
\]
and there is the exact identity
\begin{equation}
\label{eq:rowwise-covariance-centering}
 k_{N_j;g,h}-\varphi_j(k_{N_j;g,h})1
 =\frac1{N_j}\sum_{i=1}^{N_j}y_{j,i;g,h}.
\end{equation}

Let \(U_j\) be an ordered word having \(d\) fluctuation factors
and \(r\geq1\) factors of the form on the left of
Equation~\eqref{eq:rowwise-covariance-centering}. Expand all its sums.
Each of its \(d+r\) row factors is centered under
\(\varphi_j\), and a singleton row again gives zero by
Lemma~\ref{lem:factor-independence}. A surviving term therefore
uses at most \(\lfloor(d+r)/2\rfloor\) distinct row labels.
The products within any one row have absolute expectation
bounded by the product of the operator norms of their
factors. There are only finitely many partitions of the
\(d+r\) positions. Together with the normalization
\(N_j^{-d/2-r}\), these observations give
\begin{equation}
\label{eq:centered-mixed-word-bound}
 \left|\varphi_j(U_j)\right|
 \leq C_{d,r,M}\,
 N_j^{\lfloor(d+r)/2\rfloor-d/2-r}
 \leq C_{d,r,M}\,N_j^{-r/2}
 \longrightarrow0,
\end{equation}
where the constant is independent of \(j\) and
\(\varphi_j\).

Finally, replace each covariance factor in a general word
by its scalar \(\varphi_j\)-mean plus the centered factor
in Equation~\eqref{eq:rowwise-covariance-centering}. Every term
containing a centered covariance factor tends to zero by
Equation~\eqref{eq:centered-mixed-word-bound}. The remaining terms
converge by the fluctuation calculation and
Equation~\eqref{eq:covariance-convergence}. This proves convergence
of every ordered word. Since all generators are
self-adjoint, linearity also gives the
\(*\)-polynomial formulation in
Equation~\eqref{eq:joint-polynomial-limit}.
\end{proof}
